\documentclass[12pt,a4paper]{article}
\usepackage[utf8]{inputenc}
\usepackage{lmodern}
\usepackage[T1]{fontenc}
\usepackage[finnish]{babel}

\usepackage{graphicx}
\usepackage{epstopdf}

\begin{document}
\title{Tietokantojen perusteet, kevät 2020 \\
Harjoitustyön raportti}
\author{Ville Manninen \\014922187 \\Macville}
\date{\today}
\clearpage\maketitle
\begin{center}
\vspace{2cm}
\includegraphics{university-of-helsinki-2.eps} 
\end{center}
\thispagestyle{empty}

\newpage
\tableofcontents
\newpage
\section{Esittely}
Harjoitustyö on toteutettu Java-ohjelmana käyttäen NetBeans ympäristöä, ohjelmalla on seuraavat ominaisuudet.
\begin{itemize}
\item Tietokannan ja puuttuvien taulukoiden luominen.
\item Asiakkaiden lisääminen tietokantaan nimellä.
\item Paikan lisäämien tietokantaan nimellä.
\item Paketin lisääminen tietokantaan käyttäen asiakkaan nimeä.
\item Tapahtuman lisääminen tietokantaan käyttäen paikannimeä ja paketin seurantakoodia.

\end{itemize}
Ohjelma tarkistaa käyttäjän antamat syötteet ja ilmoittaa syyn jos syöte on puutteellinen.

\newpage
\section{Kaaviot ja Skeemat}
\subsection*{Tietokantakaavio}
Harjoitustyön käyttämällä tietokannalla on seuraavanlainen tietokantakaavio.
\begin{center}

\includegraphics[width=1\textwidth]{TikapeKaavio.png} 
\end{center}
\subsection*{SQL-skeema}

\newpage
\section{Tehokkuustesti}
\subsection*{Ilman indeksejä}
\subsection*{Indeksien lisäämisen jälkeen}

\newpage
\section{Tietokannan eheys}
Tietokantaan ei ole mahdollista luoda samannimisiä asiakkaita tai paikkoja, eikä tietokantaan luoduilla paketeilla saa antaa samaa seurantakoodia. Samannimisiä syötteitä valvotaan SQliten toiminnolla UNIQUE, jonka avulla voidaan varmistaa että tietokantaan ei lisätä samannimisiä arvoja. Java sovellus varmistaa ettei sille annetut syötteet ole tyhjiä, joten tietokantaan ei tallenneta NULL arvoja.

Tietokantaan lisättäessä paketteja tai tapahtumia, jotka saavat arvokseen viiteavaimia, sovellus suorittaa kyselyn joka varmistaa että annetut viiteavaimet löytyvät tietokannasta.





\newpage
\section{Lähdekoodi}
\subsection*{GitHub}
\textbf{https://github.com/Viltska/sql-app-training}
\subsection*{Koodi}


\newpage

\end{document}